\documentclass[10pt]{beamer}
\usetheme[
%%% option passed to the outer theme
    %progressstyle=fixedCircCnt,   % fixedCircCnt, movingCircCnt (moving is default)
  ]{Feather}
  
% If you want to change the colors of the various elements in the theme, edit and uncomment the following lines
\definecolor{grey}{rgb}{0.52, 0.52, 0.51}

% Change the bar colors:
\setbeamercolor{Feather}{fg=grey!20,bg=grey}
\include{beamerinnerthemeFeather}
% Change the color of the structural elements:
\setbeamercolor{structure}{fg=grey}

% Change the frame title text color:
\setbeamercolor{frametitle}{fg=blue}

% Change the normal text color background:
\setbeamercolor{normal text}{fg=black,bg=gray!10}

%-------------------------------------------------------
% INCLUDE PACKAGES
%-------------------------------------------------------

\usepackage[utf8]{inputenc}
\usepackage[ngerman]{babel}
\usepackage[T1]{fontenc}
\usepackage{helvet}
\usepackage{lipsum}
\usepackage{tikz}
\usepackage[export]{adjustbox}
\usepackage{csquotes}

\usetikzlibrary{arrows, arrows.meta, positioning}

%-------------------------------------------------------
% DEFFINING AND REDEFINING COMMANDS
%-------------------------------------------------------

% colored hyperlinks
\newcommand{\chref}[2]{
  \href{#1}{{\usebeamercolor[bg]{Feather}#2}}
}

%-------------------------------------------------------
% INFORMATION IN THE TITLE PAGE
%-------------------------------------------------------

\title[] % [] is optional - is placed on the bottom of the sidebar on every slide
{ % is placed on the title page
      \textbf{Wavelength}
}

\subtitle[$\lambda$-IDE]
{
      \textbf{$\lambda$-IDE}
}

\author[wavelength]
{     
      Muhammet Gümüs \\
      Jean-Pierre von der Heydt \\  
       Markus Himmel \\
       Marc Huisinga \\
       Philip Klemens \\ 
       Julia Schmid   }

\institute[]
{
      
  
  %there must be an empty line above this line - otherwise some unwanted space is added between the university and the country (I do not know why;( )
}

\date{19. März 2018}

%-------------------------------------------------------
% THE BODY OF THE PRESENTATION
%-------------------------------------------------------

\begin{document}

%-------------------------------------------------------
% THE TITLEPAGE
%-------------------------------------------------------

%{\1 % this is the name of the PDF file for the background
% Mein PDF-Viewer bekommt das nicht hin, das Bild im Hintergrund anzuzeigen
{
\setbeamertemplate{background canvas}{\tikz[remember picture,overlay]\node[opacity=0.3] at (current page.center) {\includegraphics[height=\paperheight]{img/logo.pdf}};}
\begin{frame}[plain]
\maketitle
\end{frame}
}

\begin{frame}{JUnit}% JUnit
Komponententests der GWT unabhängigen Pakete:
\begin{itemize}
\item[•] edu.kit.wavelength.client.model
\item[•] edu.kit.wavelength.client.view.exercise
\item[•] edu.kit.wavelength.client.view.view
\item[•] edu.kit.wavelength.server.database
\end{itemize}

\end{frame}

\begin{frame}{Selenium}
Automatisierte Tests aller Pfichtenheft-Testszenarien
\end{frame}

\begin{frame} {Fuzzer}

\end{frame}

\begin{frame}{Features}
\begin{itemize}
\item Quality-of-Life Verbesserungen:
\begin{itemize}
\item[•] Tooltips
\item[•] Statusanzeiger und Reduktionsschrittzähler
\item[•] Vorschlagsfunktion für Bibliotheksterme
\item[•] Tastenkürzel zum Starten und Beenden der Auswertung
% \item[•] Workflow
\end{itemize}
\item[•] Änderung von Ausgabeformat, Ausgabeumfang und Reduktionsordnung während laufende Auswertung möglich.
\item[•] Wiederholbare Übungsaufgabe zu den Reduktionsordnungen
\end{itemize}
% wiederholbare Übungsaufgaben
% keine automatische Serialisierung
\end{frame}

\begin{frame} % Statistiken TODO Zahlen 
\begin{itemize}
\item Lines of code:
\item Anzahl Tests:
\item erstellte Issues
\item Commits in der Qualitätssicherungsphase
\end{itemize}
\end{frame}

\end{document}
